

\section*{Ollama Environment Configuration}

The Ollama model server is configured through environment variables in the Docker Compose deployment.
These settings affect runtime behavior such as model retention, cache representation, memory mapping, and optional acceleration features.
They are implementation parameters of the model-serving infrastructure and do not change the semantic extraction, evaluation, reconstruction, or vocabulary-mapping logic of the workflow.

\begin{verbatim}
OLLAMA_KV_CACHE_TYPE=q8_0
OLLAMA_KEEP_ALIVE=-1
OLLAMA_FLASH_ATTENTION=1
OLLAMA_NO_MMAP=0
\end{verbatim}

\begin{itemize}
    \item \texttt{OLLAMA\_KV\_CACHE\_TYPE=q8\_0} configures the representation of the key--value cache used during model inference.
    The prototype uses \texttt{q8\_0} to reduce memory consumption compared with an unquantized cache representation.
    \item \texttt{OLLAMA\_KEEP\_ALIVE=-1} keeps loaded models resident in memory after requests instead of unloading them between individual calls.
    This reduces repeated model-loading overhead during long workflow runs.
    \item \texttt{OLLAMA\_FLASH\_ATTENTION=1} enables Flash Attention when it is supported by the deployment environment.
    \item \texttt{OLLAMA\_NO\_MMAP=0} leaves memory mapping enabled.
\end{itemize}

In addition to these environment variables, the Docker Compose configuration can expose available NVIDIA devices to the Ollama container.
This allows the same workflow implementation to run with CPU inference, local \ac{GPU} acceleration, or a remote Ollama instance on a dedicated inference server.
